\section{Other considered ideas and claims not advanced}
\label{sec:shortlist}

These six ideas are ranked for possible reconsideration, rather than proposed as six additional full campaigns. Each has a specific unresolved decision. Existing equipment may change the order in which a small test is worthwhile.

\paragraph{1. Tungsten coordination and retention in sugar conversion: exploratory hold.}
The opportunity is to connect productive sugar coordination, W export and useful glycol output to a controllable anchoring or feed variable. Ti--O--W stabilization and reuse already have strong precedent. Li's supported-Ni control without intentionally added W is labeled Ni/M. Its 29.5\% ethylene-glycol yield with preleachate must be compared with 27.1\% for Ni/M alone; it is not a 29.5\% leachate-derived contribution.\citep{li2020w,li2020wsi} If reopened, first reproduce a retained and a soluble-W reference with hot export and carbon balances. Transferred activity can demonstrate exported function but cannot uniquely assign all upstream turnover. A chemically grounded intervention addressing a consequential reference limitation remains missing.

\paragraph{2. Product-rich liquid Ti-zeolite epoxidation: conditional model-transfer test.}
Ask whether a calibrated reaction network predicts an independent product-accumulation trajectory, including peroxide loss and epoxide ring opening. Product-induced solvent changes already have direct spectroscopic precedent, and related Ti-zeolite secondary kinetics contain unresolved behavior.\citep{torres2023thesis,kwon2024} Qualify component kinetics in the chosen liquid catalyst/solvent system before freezing the prediction; a vapor-fed model is not a validated liquid baseline. Close the additional-mechanism claim if the established network predicts useful output adequately. A residual alone does not identify solvent reorganization, fouling or Ti loss.

\paragraph{3. Zirconia styrene self-cleaning: bounded chemical-balance pilot.}
Hydrocarbon-assisted removal of inhibiting water is already proposed for Zr--O acid--base pairs; the cleaning products were not identified in the source's dilute-feed regime.\citep{artsiusheuski2025} First establish useful net output with measured moisture and relevant styrene/hydrogen burdens, then cross water history with ethylbenzene exposure to relate oxygen export to recovered function. Control oxygen exchange and perturbation by the activity assay before isotope attribution. Expand only if the chemistry predicts a useful feed or removal policy; omitting steam alone does not establish process savings.

\paragraph{4. Does an acid-site assay change subsequent zeolite aging? Conditional protocol study.}
Compare uninterrupted aging of protonic chabazite (H-CHA) with repeated NH\textsubscript{3}/water assays and a matched thermal/water sham. Measure function before terminal titration on separate matched specimens; a common terminal assay then compares recoverable capacity. Repeated assays are part of the relevant published protocol; a Cu/SSZ-13 aging-order comparison also assayed both arms and therefore does not answer assay versus no assay.\citep{classmartinez2026,luo2018} The useful result would predict later function under a withheld assay spacing. Immediate uptake recovery is insufficient, and a null schedule effect cannot by itself determine the lifetime of an unseen Al intermediate.

\paragraph{5. Wet methane oxidation after NO removal: application-dependent diagnostic.}
Test whether the sulfur-tolerance advantage of Pd/SSZ-13 treated by ST-CO-N2-O2 over its ST-CO-O2 counterpart survives removing NO at matched water, sulfur history, temperature and precious-metal inventory. The former includes an extra N\textsubscript{2} treatment between CO and O\textsubscript{2}. The reported sulfur interval includes NO, while history-dependent NO promotion on other Pd catalysts is already known.\citep{ryu2024,auvinen2021} Apply the small crossed NO/sulfur test to both variants only where a real low-NO, sulfur-containing duty could change the material choice. Without that application boundary, another NO response is a diagnostic rather than a substantial program.

\paragraph{6. Ga-based propane dehydrogenation regeneration: conditional cycle study.}
Test whether a measured restoration endpoint predicts a shorter useful regeneration across a withheld spent history. Redox-related deactivation and shortened oxygen treatments are already reported.\citep{wu2025,buchbinder2022} Oxygen consumption minus \COtwo{} alone is not a selective Ga-redox inventory: oxidation of retained hydrogen and other oxygen-containing products also enters the balance. Calibrate that balance before assigning sites, and compare complete-cycle propylene output, gas use and heat duty with the best simple regeneration policy. Correcting an assay does not by itself establish a better catalyst cycle.

\subsection*{Claims withdrawn, rather than topics declared invalid}

Several attractive shortcuts do not survive scrutiny. A faster outlet-water transient does not uniquely establish lower active-site exposure. Fresh-Na rescue does not identify selective initiation loss. Ni compatibility and retention on promoted Ag are already prior art. The detailed chapters replace these claims with narrower functional and predictive questions; none of these corrections invalidates its entire catalytic field.

For tungsten, the inactive-mannose-reservoir premise should also be withdrawn: published hot mannose feeds produce glycolaldehyde, with epimerization complicating assignment of the direct cleavage route.\citep{elmohammad2024} This does not exclude an inactive fraction, but mannose cannot be assumed to form an inactive reservoir. The proposal remains on hold until a better supported intervention is selected. More broadly, the shortlist distinguishes missing evidence from a demonstrated contradiction: an incomplete program can be reconsidered after its specific missing test, whereas an invalid inference should not be retained as the default mechanism.
